% Audit du code des thématiques à organismes proposés de mod_stage : constats, corrections, tests.
% Compilation : latexmk -pdf audit-organismes-proposes.tex
\documentclass[11pt,a4paper]{article}
\usepackage[french]{babel}
\usepackage[T1]{fontenc}
\usepackage{lmodern,microtype,booktabs,tabularx,xcolor,hyperref,enumitem}
\usepackage[margin=2.2cm]{geometry}
\usepackage{url}
\expandafter\def\expandafter\UrlBreaks\expandafter{\UrlBreaks\do\a\do\b\do\c\do\d\do\e\do\f\do\g\do\h\do\i\do\j\do\k\do\l\do\m\do\n\do\o\do\p\do\q\do\r\do\s\do\t\do\u\do\v\do\w\do\x\do\y\do\z}
\DeclareUrlCommand\code{\urlstyle{tt}}
\Urlmuskip=0mu plus 2mu
\newcommand{\grav}[1]{\textbf{\textcolor{#1}{\textbullet}}}
\newcommand{\maj}{\grav{orange}}
\newcommand{\mineure}{\grav{gray}}
\newcommand{\crit}{\grav{red}}
\setlist[itemize]{itemsep=2pt,topsep=3pt}
\newcommand{\ok}{\textcolor{teal}{\textbf{corrigé}}}
\newcommand{\partiel}{\textcolor{orange}{\textbf{en partie}}}
\newcommand{\non}{\textcolor{gray}{\textbf{non traité}}}
\title{Audit des thématiques à organismes proposés de \texttt{mod\_stage}}
\author{Sébastien Lefebvre}
\date{17 novembre 2026}

\begin{document}
\maketitle

\section{Objet et méthode}

Ce rapport couvre le code ajouté depuis le commit \code{29099d0} sur la branche
\code{ccr-9819324d-uc8w8o} : les thématiques à organismes proposés (commit \code{341c6a7}), leurs
tests et leur test de charge (commits \code{6690f9a} et \code{82fce02}). Il porte sur
\code{placementlib.php}, les six pages \code{placement*.php}, le formulaire et l'importateur
d'organismes, la tâche planifiée, et les parties modifiées de la demande de convention, de la fiche
de thématique, de la sauvegarde, de la restauration et du fournisseur de données personnelles.
Le module est en version \code{2026111700}.

\begin{itemize}
  \item \textbf{Une revue indépendante}, conduite sans connaître les conclusions de l'auteur, sur
    sept axes : sécurité, contournement des règles, concurrence, portabilité SQL, sauvegarde et
    données personnelles, performance, code mort et chaînes de langue. Chaque constat est assorti
    d'un scénario d'échec et a été vérifié à la lecture de toutes les lignes en cause.
  \item \textbf{La relecture de l'auteur}, menée en parallèle. Trois de ses constats recoupent ceux
    de la revue ; quatre s'y ajoutent.
  \item \textbf{Les tests joués}, contrairement à l'audit précédent : PHPUnit sur un Moodle~4.5
    et PostgreSQL~16, le contrôle de style \code{phpcs} avec le standard \code{moodle}, et le banc
    d'utilisateurs simulés de \code{docs/audit/organismes/}, rejoué après les corrections.
  \item \textbf{MariaDB et Behat}, joués dans un second temps (section~\ref{sec:verifications}).
  \item \textbf{Limites} : Moodle~4.3 avec PHP~8.0, borne basse de la CI, n'a pas été joué ici. La
    tâche d'envoi des courriels a été exécutée par les tests, pas par un cron réel.
\end{itemize}

\noindent\textbf{Échelle} : celle des rapports précédents. \crit~critique, \maj~majeure,
\mineure~mineure.

\section{Synthèse}

L'audit n'a trouvé \textbf{aucun défaut critique et aucune faille de sécurité}. Il a relevé
\textbf{trois défauts majeurs, seize mineurs et deux points de finition}. Tous sont corrigés, sauf
un point de finition traité en partie. Une quatrième anomalie majeure est née d'une correction de
cet audit, et a été rattrapée par le test de charge avant d'être livrée
(section~\ref{sec:regression}).

\begin{center}
\small
\begin{tabularx}{\linewidth}{@{}l X l@{}}
\toprule
& \textbf{Constat} & \textbf{État} \\
\midrule
\maj & Règle du SIRET contournable en cochant \og à l'étranger\fg{} sans indiquer de pays & \ok \\
\maj & Après une réinitialisation du cours, la campagne de vœux se refermait seule dans l'heure & \ok \\
\maj & L'étudiant affecté pouvait déposer sa demande pour un autre organisme que le sien & \ok \\
\maj & Import de 2\,000 lignes en échec, par une correction de cet audit & \ok \\
\mineure & Affectation déplacée ou retirée alors que la demande de convention est déposée & \ok \\
\mineure & Affectation libre non relue sous le verrou en mode édition : deux demandes pour une affectation & \ok \\
\mineure & Demande d'un stage déjà enregistré : affectation de l'étudiant ignorée, puis réutilisable & \ok \\
\mineure & Affectation consommée par un stage enregistré chez un autre organisme & \ok \\
\mineure & Tâche planifiée : appariement d'une thématique dont la date vient d'être repoussée & \ok \\
\mineure & Tâche planifiée : une thématique en échec privait les suivantes de leur appariement & \ok \\
\mineure & Thématique masquée après l'appariement : plus aucun accès à la demande préremplie & \ok \\
\mineure & Organismes de même nom et de SIRET différents confondus à l'import et à la reprise & \ok \\
\mineure & Suppression et reprise d'organismes hors verrou : affectation orpheline, doublons & \ok \\
\mineure & Date de clôture passée acceptée, y compris à la réouverture des vœux & \ok \\
\mineure & Pays \og France\fg{} importé : organisme traité comme étranger & \ok \\
\mineure & Courriels de résultat envoyés dans la requête de la page & \ok \\
\mineure & Échange de deux étudiants impossible, malgré le commentaire du code & \ok \\
\mineure & Organisme enregistrable sans nom, ou avec moins de places que d'affectés & \ok \\
\mineure & Noms sans lettre ni chiffre tous confondus entre eux & \ok \\
\mineure & Nom de l'étudiant non échappé dans un message d'erreur & \ok \\
\mineure & Test unitaire fondé sur une durée d'exécution & \ok \\
& Export des données personnelles : rien pour le membre de la DEVE cité & \ok \\
& Quelques requêtes par thématique dans l'encart de l'étudiant et l'export personnel & \partiel \\
\bottomrule
\end{tabularx}
\end{center}

\section{Défauts majeurs}

\begin{itemize}
  \item[\maj] \textbf{SIRET contournable.} La règle du hors liste dispensait de SIRET dès que la
    case \og à l'étranger\fg{} était cochée, sans regarder le pays. Un étudiant cochait la case,
    laissait le pays vide et déposait sa demande sans SIRET ; le stage comptait en plus comme stage
    à l'étranger. \emph{Correction} : la dispense exige un pays, qui ne soit pas la France
    (\code{stage_placement_request_errors()}). Sans pays, le formulaire demande le pays ou le SIRET.
  \item[\maj] \textbf{Campagne refermée après réinitialisation.} La réinitialisation du cours
    effaçait la date de l'appariement mais gardait la date de clôture, déjà passée. La thématique
    redevenait \og vœux clos\fg{}, la tâche planifiée l'appariait sans aucun vœu dans l'heure, et
    la nouvelle promotion trouvait une campagne terminée. Même effet après une suppression des
    données personnelles de toute l'activité. \emph{Correction} : la date de clôture est effacée
    elle aussi, comme le fait déjà une restauration sans les étudiants.
  \item[\maj] \textbf{Demande pour un autre organisme.} Une demande ouverte à partir d'une
    affectation n'était plus contrôlée du tout : l'étudiant pouvait remplacer l'organisme attribué
    par un autre, sans SIRET, même sur une thématique où le hors liste est interdit. L'affectation
    était marquée utilisée et la place restait prise. \emph{Correction} : la demande doit porter
    sur l'organisme attribué (même nom, et même SIRET s'il en a un). Les autres coordonnées restent
    corrigeables. C'est un choix : voir la section~\ref{sec:decision}.
\end{itemize}

\section{Régression introduite et rattrapée}\label{sec:regression}

\begin{itemize}
  \item[\maj] \textbf{Import de 2\,000 lignes en échec.} Pour ne plus confondre deux organismes de
    même nom, la première correction cherchait chaque ligne du classeur en parcourant tous les
    organismes connus, et relisait l'organisme en base après chaque ligne. Le coût devenait
    quadratique. Les 46~tests unitaires passaient ; le banc d'utilisateurs simulés, lui, a renvoyé
    deux erreurs~500 sur l'import et le réimport de 2\,000~lignes, et un débit divisé par deux.
    \emph{Correction} : un index par SIRET et par nom
    (\code{stage_placement_index_offers()}). L'import de 2\,000~lignes prend de nouveau deux secondes
    environ.
\end{itemize}

Cette anomalie n'a jamais été poussée. Elle montre ce que le banc apporte que les tests unitaires
n'apportent pas : ceux-ci vérifiaient le résultat sur quelques lignes, pas le coût sur un volume.

\section{Défauts mineurs}

\subsection{Affectations}

\begin{itemize}
  \item[\mineure] \textbf{Affectation modifiée après le dépôt.} La DEVE pouvait déplacer ou
    retirer un étudiant dont la convention était déjà demandée : le tableau et la convention
    disaient alors deux organismes différents. \emph{Correction} : refus, avec un message qui
    demande d'annuler d'abord le stage.
  \item[\mineure] \textbf{Deux demandes pour une affectation.} En mode édition d'un stage
    enregistré par la DEVE, l'affectation libre de l'étudiant n'était pas relue sous le verrou. Avec
    deux onglets, elle servait à un nouveau stage puis était rattachée à l'ancien.
    \emph{Correction} : relecture sous le verrou, et refus si elle a servi entre-temps.
  \item[\mineure] \textbf{Page \code{convention_request.php}.} Encore accessible à l'étudiant par
    son adresse, elle ignorait son affectation libre : elle exigeait le SIRET, puis laissait
    l'affectation réutilisable. \emph{Correction} : même règle que l'autre formulaire, et
    rattachement du stage à l'affectation.
  \item[\mineure] \textbf{Affectation consommée à tort.} Un stage enregistré par la DEVE chez un
    autre organisme était rattaché à l'affectation libre de l'étudiant. \emph{Correction} :
    l'affectation n'est reprise que si le stage est chez l'organisme attribué
    (\code{stage_placement_assignment_for_entry()}) ; sinon la demande est hors liste.
  \item[\mineure] \textbf{Échange de deux étudiants.} Les affectations du tableau étaient
    enregistrées une à une : deux étudiants ne pouvaient pas échanger leurs organismes complets.
    \emph{Correction} : le lot est appliqué en entier ou pas du tout, les places libérées par le
    lot étant disponibles pour lui (\code{stage_placement_set_assignments()}).
\end{itemize}

\subsection{Appariement et campagne}

\begin{itemize}
  \item[\mineure] \textbf{Thématique périmée.} La tâche planifiée lit la liste des thématiques une
    fois, puis envoie des courriels pendant plusieurs minutes. Une date repoussée entre-temps
    n'empêchait pas l'appariement. \emph{Correction} : la thématique est relue sous le verrou, et
    l'appariement automatique n'a lieu que si les vœux sont toujours clos.
  \item[\mineure] \textbf{Échec en chaîne.} Une thématique dont l'activité était verrouillée
    interrompait la tâche. \emph{Correction} : l'échec est consigné, les thématiques suivantes
    sont traitées, la thématique en échec est reprise au passage suivant.
  \item[\mineure] \textbf{Date de clôture passée.} Elle était acceptée à l'activation et à la
    réouverture des vœux, qui restait alors sans effet. \emph{Correction} : refus quand la date
    vient d'être saisie ou que les vœux sont rouverts. Une date passée déjà enregistrée ne bloque
    pas le reste de la fiche.
  \item[\mineure] \textbf{Thématique masquée.} L'étudiant affecté perdait tout lien vers sa
    demande, et le courriel de résultat menait à une page indisponible. \emph{Correction} :
    l'encart et la page des organismes restent ouverts à l'étudiant qui a des vœux ou une
    affectation.
  \item[\mineure] \textbf{Courriels dans la page.} Lancé par la DEVE, l'appariement envoyait un
    courriel par étudiant avant de répondre. Pour une promotion entière, la page pouvait dépasser
    son délai, et les étudiants non prévenus ne pouvaient plus l'être. \emph{Correction} : l'envoi
    est confié à une tâche ponctuelle (\code{notify_placement_results}), exécutée par le cron.
\end{itemize}

\subsection{Organismes}

\begin{itemize}
  \item[\mineure] \textbf{Même nom, SIRET différents.} Deux sites d'une même clinique, chacun avec
    son SIRET, étaient confondus : la seconde ligne du classeur écrasait la première.
    \emph{Correction} : le nom ne rapproche deux organismes que si l'un des deux n'a pas de SIRET.
  \item[\mineure] \textbf{Écritures hors verrou.} La suppression d'un organisme pendant un
    appariement pouvait laisser une affectation orpheline ; un double clic sur la reprise de
    stages précédents créait des doublons. \emph{Correction} : verrou de l'activité et
    transaction dans les deux fonctions.
  \item[\mineure] \textbf{Pays \og France\fg{}.} Une colonne \og Pays\fg{} remplie de
    \og France\fg{} faisait de chaque organisme un organisme étranger. \emph{Correction} : ce pays
    est effacé à l'enregistrement.
  \item[\mineure] Organisme sans nom (nom fait d'espaces) ; places ramenées sous le nombre
    d'étudiants affectés ; noms sans lettre ni chiffre confondus entre eux ; nom de l'étudiant non
    échappé dans le message d'erreur des affectations ; test unitaire qui échouait quand la
    machine ralentissait. \emph{Tous corrigés}.
\end{itemize}

\section{Finition}

\begin{itemize}
  \item L'export des données personnelles ne sortait rien pour le membre de la DEVE qui avait
    décidé des affectations. Il exporte maintenant ses décisions : thématique, organisme et date,
    sans l'identité de l'étudiant. \ok
  \item L'encart de l'étudiant faisait jusqu'à quatre requêtes par thématique. Les affectations
    sont maintenant lues en une requête ; il reste une requête de comptage par thématique dont les
    vœux sont ouverts. L'export des données personnelles fait toujours trois requêtes par
    thématique. Sans effet mesurable tant que les thématiques se comptent en dizaines. \partiel
\end{itemize}

\section{Décision à confirmer}\label{sec:decision}

La correction du troisième défaut majeur impose que la demande ouverte à partir d'une affectation
porte sur l'organisme attribué. C'est la lecture la plus stricte de la fonction. L'autre lecture
est possible : laisser l'étudiant changer d'organisme et le signaler à la DEVE lors de l'instruction
de la convention. Elle serait plus souple, mais une thématique où le hors liste est interdit
pourrait alors être contournée par tout étudiant affecté. Si la souplesse est préférée, la règle
se desserre en un point : \code{stage_placement_request_errors()}.

\section{Ce qui tient}

\begin{itemize}
  \item \textbf{Sécurité.} Chaque page vérifie la connexion et la capacité, charge la thématique
    avec son activité et l'organisme avec sa thématique. Toutes les écritures sont en POST avec la
    clé de session. La reprise de stages précédents ne montre que les activités où l'utilisateur
    gère les thématiques. Les sorties sont échappées. Aucune redirection ouverte.
  \item \textbf{Contournements testés et bloqués.} Thématique d'une autre activité dans
    l'adresse, changement de thématique en mode édition, affectation utilisée deux fois par double
    envoi, vœux après la clôture ou après l'appariement.
  \item \textbf{Verrou.} Aucun chemin ne prend le verrou de l'activité deux fois. Il est libéré
    dans tous les cas, y compris sur les sorties anticipées, après la validation de la transaction.
  \item \textbf{Portabilité SQL.} Pas de sous-requête sur la table modifiée, pas de mot réservé.
    \code{install.xml} et \code{upgrade.php} sont identiques champ par champ ; la mise à jour
    depuis la version précédente, jouée, donne le schéma d'une installation neuve.
  \item \textbf{Sauvegarde et données personnelles.} Vœux et affectations ne suivent que les
    données utilisateur ; les identifiants sont mis en correspondance, y compris le stage rattaché.
  \item \textbf{Chaînes de langue.} Toutes les chaînes utilisées sont définies en français et en
    anglais, et toutes celles définies sont utilisées.
\end{itemize}

\section{Vérifications après correction}\label{sec:verifications}

\begin{itemize}
  \item \textbf{PHPUnit} : 46~tests pour la fonction, dont un par correction de cet audit. La suite
    de \code{mod_stage} et celle de \code{mod_stagesynthesis} passent sur Moodle~4.5 et PostgreSQL.
  \item \textbf{MariaDB~10.11} : les 333~tests de \code{mod_stage} et les 8 de
    \code{mod_stagesynthesis} passent aussi sur un second Moodle branché sur MariaDB.
  \item \textbf{Behat} : quatre scénarios ajoutés (\code{tests/behat/placement.feature}) : vœux,
    appariement et demande préremplie ; hors liste refusé ; hors liste admis avec SIRET ou pays ;
    gestion des organismes et affectation à la main par la DEVE. Les 15~scénarios du module
    passent, sans JavaScript, comme dans la CI.
  \item \textbf{Style} : \code{phpcs} propre sur tout le dépôt.
  \item \textbf{Banc d'utilisateurs simulés} : les quatre vagues ont été rejouées sur le code
    corrigé. Leurs résultats sont dans \code{test-de-charge-organismes.tex} : aucune erreur, neuf
    invariants respectés. Une cinquième vague, sur MariaDB, donne le même résultat.
\end{itemize}

\section{Suite}

\begin{itemize}
  \item Faire jouer la CI avant toute fusion, pour la borne basse (Moodle~4.3 avec PHP~8.0).
  \item Confirmer la décision de la section~\ref{sec:decision}.
\end{itemize}

\end{document}
